\appendix
\section{复算路径与核心伪代码}
本论文绑定 model v42、code v15、solve v12。结构化主结果为 \texttt{data/results.json}，Q1 至 Q4 的明细分别位于 \texttt{q1\_forecast\_metrics.csv}、\texttt{q2\_named\_scheme\_comparison.csv}、\texttt{q3\_search\_trace.csv} 和 \texttt{q4\_scenario\_comparison\_v42.csv}。图形只读取同一 solve v12 的 CSV 或其正式图形清单。


\begin{table}[H]\centering
{\small
\caption{正文结论与结构化证据的对应关系}
\begin{tabular}{lll}\toprule
问题 & 正文核心结论 & 复算证据\\\midrule
Q1 & 闭环预测与 538 任务末端调度 & q1\_forecast\_metrics、task\_schedule\\
Q2 & 五方向结果与购售电分解 & q2\_named\_scheme、q2\_energy\\
Q3 & 七策略、六轨迹及循环次数 & q3\_search\_trace、q3\_storage\_energy\\
Q4 & 动态交替及 17 个情景 & q4\_alternation\_trace、q4\_scenario\\\bottomrule
\end{tabular}
}
\end{table}

复算时先核对版本绑定与文件哈希，再检查结构化结果是否齐全，最后执行约束审计。数值通过不等于方法通过：若结果文件存在但方法合同中的目标、约束、算法类别或近似声明不一致，论文仍不得沿用该数值。同样，约束通过不等于最优；终止状态为候选预算耗尽、最大迭代或连续无改善时，只能报告当时的可行 incumbent。

图表也采用相同边界。每幅数值图必须能追溯到 solve v12 当前 CSV；旧 paper 构建目录中遗留但未出现在当前图形清单的图片不进入 v20。预测比较图由当前预测明细按小时聚合，Q1 GPU、Q3 SOC 及两幅灵敏度图直接来自 solve v12 正式图形。论文中的亿元、万元和百分比只是展示换算，审计仍以原始元、MWh、MW 与无量纲值为准。

\begin{verbatim}
read_and_validate_inputs()
fit Q1 candidates on train; select on validation
closed_loop_test(); schedule actual last-24h tasks
for direction in five_Q2_directions:
    evaluate 10000 incremental task moves
    audit all 50000 tasks and hourly energy
for policy in seven_Q3_policies:
    run 21-state SOC dynamic program
    audit, hash, deduplicate, Pareto-filter
for scenario in seventeen_Q4_scenarios:
    restart from common schedule and initial SOC
    repeat up to 8 rounds:
        optimize storage for current task load
        rebuild 600 task moves from current energy signal
        accept only audited strict improvement
        stop after 2 consecutive non-improving rounds
write atomically; run final constraint audit
\end{verbatim}

方法的正式实现为 code v15 的 \texttt{solution.py}。本次 v20 只重建论文与审核证据，没有修改该程序、solve v12 结构化结果或既有 v17/v18 文件。
